\section{From squares to cubes}\label{sec:story}
The method was developed on unit squares in a square, where records are dense and audited \citep{squaresproject,friedmansurvey}. With the original schedule, hardening reaches Trump's packing for $n=11$ \citep{trump} after tightening (to within $10^{-14}$) in 4 of 80 random starts at three times the base budget and 2 of 80 at the base budget; rigid starts with the same schedule reach it in none of 80 at either budget (campaign C00).

An earlier note by the author, \emph{Twelve unit cubes in a cube of side 2.93152} (submitted to Geombinatorics; current editorial status unconfirmed), reported a feasible packing of 12 unit cubes in a cube of side $2.9315185$. The public code and certificate are separately available \citep{softtorigidpacking}. Its public repository compares this bound with the catalogued Hyra value $2.9327718$ and Friedman's 1998 value $2.94281$ \citep{hyra,friedmancubes}; these are historical comparator snapshots, not a current priority adjudication. The author's public build log flags novelty relative to Haowei Lin's arrangement as disputed \citep{prioritycontext2026}. We therefore use the earlier bound only as context and make no new-record or undisputed-priority claim; its priority is unnecessary to the controlled comparison here. Campaign C05 repeats the cube experiment with three gradient paths and two settings at the earlier bound's search budget (Table~\ref{tab:cubes}). Hardening has the lowest side of the three paths at $n=11$, 12 and 13 under both settings, all paths reach Friedman's $n=9$ packing \citep{friedmancubes}, and with tuned settings hardening and rigid starts also reach his $n=10$ packing $2+1/\sqrt2$, which the original schedule never did. Rigid starts end at the trivial side 3 for $n\ge12$. In these 32 seeds nothing went below a catalogue value: the best side at $n=12$ is \CubesHardenTwelve{} (grow \CubesGrowTwelve, rigid \CubesRigidTwelve), while the earlier feasible bound was reported from one of 192 runs. At $n=13$ hardening reaches $2.956145$, consistent with Friedman's $2.956+$ \citep{friedmancubes} to the precision printed.

\begin{table}[t]\centering\small
\begin{tabular}{rrrrrrrr}
\toprule
 & & \multicolumn{3}{c}{original settings} & \multicolumn{3}{c}{2D-tuned settings}\\
$n$ & reference & harden & grow & rigid & harden & grow & rigid\\
\midrule
9 & 2.70711 & \textbf{2.70711} & \textbf{2.70711} & \textbf{2.70711} & \textbf{2.70711} & \textbf{2.70711} & \textbf{2.70711}\\
10 & 2.70711 & 2.76777 & 2.76777 & 2.89443 & \textbf{2.70711} & 2.76777 & \textbf{2.70711}\\
11 & 2.88295 & \textbf{2.89443} & 2.91739 & 2.91899 & 2.91250 & 2.91328 & 2.94281\\
12 & 2.93277 & 2.94301 & 2.98587 & 3.00000 & \textbf{2.94281} & 2.94751 & 3.00000\\
13 & 2.95600 & 2.99497 & 3.00000 & 3.00000 & \textbf{2.95615} & 3.00000 & 3.00000\\
14 & 2.98995 & \textbf{3.00000} & \textbf{3.00000} & \textbf{3.00000} & \textbf{3.00000} & \textbf{3.00000} & \textbf{3.00000}\\
\bottomrule
\end{tabular}
\caption{Cubes in a cube (C05): lowest side over 32 runs at the prior search budget, after tightening. The reference column excludes the prior reported $n=12$ feasible bound ($2.93152$), from a separate 192-start search; this is not a current-study catalogue target. Bold: lowest in the row.}
\label{tab:cubes}
\end{table}


Two results made the broader question worth asking. Rigid starts in these first experiments used the hardening schedule untouched, so they were never given a fair chance. And squares in a square may be special: many of its best packings tilt a few squares against their neighbours, exactly the configurations a late commitment to orientation should help with.
